% =====================================================================
%  Slide: ETSI GS QKD 018, orchestrazione SDN di reti QKD multi-dominio
%  ---------------------------------------------------------------------
%  COMPILAZIONE (pdfLaTeX normale; anche XeLaTeX/LuaLaTeX, con Cambria Math):
%      latexmk -pdf main.tex
%  oppure, a mano:
%      pdflatex main.tex ; biber main ; pdflatex main.tex ; pdflatex main.tex
%  Le figure TikZ in figures/ sono quelle del paper (paper/it/figures),
%  senza l'ambiente figure, ridimensionate con \slidefig.
% =====================================================================
\documentclass[10pt]{beamer}
% =====================================================================
%  Preambolo delle slide: tema Beamer classico (Warsaw), font, listati, TikZ
% =====================================================================

% ---------- Font e lingua / fonts and language ----------
% pdfLaTeX: Latin Modern standard. XeLaTeX/LuaLaTeX: Cambria Math se presente.
\usepackage{iftex}
\ifPDFTeX
  \usepackage[utf8]{inputenc}
  \usepackage[T1]{fontenc}
  \usepackage{lmodern}
  \usepackage{amsmath,amssymb}
\else
  \usefonttheme{professionalfonts}
  \usepackage{fontspec}
  \usepackage{amsmath,amssymb}
  \usepackage{unicode-math}
  \IfFontExistsTF{Cambria Math}
    {\setmathfont{Cambria Math}}
    {\setmathfont{Latin Modern Math}}
\fi
\usepackage[italian]{babel}
\usepackage{csquotes}

% ---------- Tema classico ----------
\usetheme{Warsaw}
\setbeamertemplate{navigation symbols}{}
% Piè di pagina diviso come nel tema, con il numero di slide a destra
\setbeamertemplate{footline}{%
  \leavevmode%
  \hbox{%
  \begin{beamercolorbox}[wd=.5\paperwidth,ht=2.5ex,dp=1.125ex,leftskip=.3cm plus1fill,rightskip=.3cm]{author in head/foot}%
    \usebeamerfont{author in head/foot}\insertshortauthor
  \end{beamercolorbox}%
  \begin{beamercolorbox}[wd=.5\paperwidth,ht=2.5ex,dp=1.125ex,leftskip=.3cm,rightskip=.3cm]{title in head/foot}%
    \usebeamerfont{title in head/foot}\insertshorttitle\hfill\insertframenumber{} / \inserttotalframenumber
  \end{beamercolorbox}}%
  \vskip0pt}

% Indice all'inizio di ogni sezione
\AtBeginSection[]{%
  \begin{frame}<beamer>
    \tableofcontents[currentsection,hideothersubsections]
  \end{frame}}

% ---------- Pacchetti ----------
\usepackage{booktabs,tabularx,array}
\usepackage{listings}
\usepackage{adjustbox}
\usepackage{tikz}
\usetikzlibrary{arrows.meta,positioning,fit,calc,shapes.geometric,backgrounds}
\usepackage[backend=biber,style=numeric-comp,sorting=none,maxbibnames=3,
            url=false,doi=false,eprint=false,isbn=false]{biblatex}
\addbibresource{bibliography/references.bib}
\renewcommand*{\bibfont}{\scriptsize}
\setbeamertemplate{bibliography item}{\insertbiblabel}
\newcolumntype{L}{>{\raggedright\arraybackslash}X}

% ---------- Listati (YANG, alberi pyang, HTTP) ----------
\lstdefinelanguage{yang}{
  morekeywords={module,namespace,prefix,import,organization,contact,
    description,revision,reference,yang-version,identity,base,typedef,type,
    identityref,grouping,uses,container,list,key,leaf,leaf-list,config,
    augment,rpc,input,output,notification,leafref,path,default,units,
    must,when,enum,enumeration,mandatory},
  sensitive=true,
  morecomment=[l]{//},
  morecomment=[s]{/*}{*/},
  morestring=[b]"
}
\lstset{
  basicstyle=\ttfamily\scriptsize,
  keywordstyle=\bfseries,
  frame=single,
  columns=fullflexible, keepspaces=true,
  breaklines=true, upquote=true,
  aboveskip=4pt, belowskip=4pt
}

% ---------- Stili TikZ comuni (come nel paper) ----------
\tikzset{
  >=Stealth,
  box/.style={draw,rounded corners,minimum height=8mm,align=center,
              font=\small},
  lbl/.style={font=\scriptsize,midway,fill=white,inner sep=1pt},
  actor/.style={draw,rounded corners,fill=gray!8,minimum width=2.4cm,
                minimum height=7mm,align=center,font=\small},
  msg/.style={->,thick},
  ret/.style={->,dashed}
}

% Figura TikZ del paper adattata alla slide: \slidefig[altezza max]{file}
\newcommand{\slidefig}[2][0.75\textheight]{%
  \begin{center}
    \begin{adjustbox}{max width=\textwidth,max totalheight=#1}
      \input{figures/#2}
    \end{adjustbox}
  \end{center}}

% Fonte in basso a destra
\newcommand{\source}[1]{\vfill\hfill{\tiny #1}}


\title[ETSI GS QKD 018]{Orchestrazione SDN di reti QKD multi-dominio}
\subtitle{L'interfaccia ETSI GS QKD 018}
\author[D.~Carro, F.~E.~Bucci]{Diego Carro \and Bucci Felice Elia}
\institute[PoliTo]{Corso di Cybersecurity, Politecnico di Torino\\Quantum Engineering}
\date{Novembre 2026}

\begin{document}

% ---------------------------------------------------------------------
\begin{frame}
  \titlepage
\end{frame}

\begin{frame}{Indice}
  \tableofcontents[hideallsubsections]
\end{frame}

% =====================================================================
\section[Introduzione]{Introduzione: dalla QKD alla rete}
% =====================================================================

\subsection{Motivazione}

\begin{frame}{Il problema}
  \begin{itemize}
    \item La QKD permette a due parti di concordare una chiave segreta con
          sicurezza fondata sulla \alert{fisica}, non sulla difficoltà
          computazionale.
    \item Portarla in una rete di operatore significa gestirla come
          qualunque altra \alert{risorsa di rete}:
          \begin{itemize}
            \item scoprirne la topologia e monitorarne lo stato;
            \item allocarne la capacità ai servizi;
            \item coordinarla con la rete ottica classica.
          \end{itemize}
    \item L'ETSI ISG QKD ha affrontato il problema per gradi:
  \end{itemize}
  \vspace{1ex}
  \begin{center}
  \begin{tikzpicture}[node distance=4mm,
      st/.style={draw,rounded corners,fill=blue!5,align=center,
                 font=\footnotesize,minimum height=11mm,text width=2.75cm}]
    \node[st] (a) {\textbf{GS QKD 002}\\casi d'uso e requisiti};
    \node[st,right=of a] (b) {\textbf{GS QKD 015}\\controller SDN $\leftrightarrow$ nodi};
    \node[st,right=of b,fill=orange!15] (c) {\textbf{GS QKD 018}\\orchestrator $\leftrightarrow$ controller};
    \draw[->,thick] (a) -- (b);
    \draw[->,thick] (b) -- (c);
  \end{tikzpicture}
  \end{center}
\end{frame}

\begin{frame}{Stato dell'arte in breve}
  \small
  \begin{columns}[T]
    \begin{column}{0.5\textwidth}
      \begin{block}{Fondamenti e reti}
        \begin{itemize}
          \item Review dei protocolli e delle prove di sicurezza~\cite{pirandola2020}
          \item Dal link alla rete: strati, trusted relay,
                instradamento~\cite{mehic2020,cao2022,routing2024}
        \end{itemize}
      \end{block}
      \begin{block}{Standard}
        \begin{itemize}
          \item ETSI: casi d'uso~\cite{qkd002}, API~\cite{qkd004,qkd014},
                SDN~\cite{qkd015}, orchestrazione~\cite{qkd018}, KMS~\cite{qkd020}
          \item ITU-T Y.3800: architettura, controllo, SDN~\cite{ituy3800,ituy3804,ituy3805}
        \end{itemize}
      \end{block}
    \end{column}
    \begin{column}{0.5\textwidth}
      \begin{block}{SDN e orchestrazione}
        \begin{itemize}
          \item Primi esperimenti SDN + QKD (UPM/Telefónica)~\cite{aguado2016,aguado2020}
          \item KOREN multi-vendor (SK Telecom)~\cite{sim2023}
          \item Prototipo gerarchico ONOS (PoliMi)~\cite{vicario2025}
          \item MadQCI, Madrid~\cite{madqci2024}
        \end{itemize}
      \end{block}
      \begin{exampleblock}{Contributo}
        Lettura unitaria della catena 002 $\to$ 015 $\to$ 018: dai
        requisiti ai modelli YANG all'orchestrazione multi-dominio.
      \end{exampleblock}
    \end{column}
  \end{columns}
\end{frame}

\subsection{QKD e trusted node}

\begin{frame}{La distribuzione quantistica di chiavi}
  \begin{columns}[T]
    \begin{column}{0.56\textwidth}
      \textbf{BB84 in quattro passi}
      \begin{enumerate}
        \item Alice codifica bit casuali su singoli fotoni, in una di due
              basi scelte a caso.
        \item Bob misura in una base casuale.
        \item Su un canale classico \alert{autenticato} confrontano le basi
              (\emph{sifting}).
        \item Correzione d'errore e \emph{privacy amplification}
              distillano la chiave.
      \end{enumerate}
      \vspace{1ex}
      Un intercettatore introduce errori misurabili:
      \[
        \mathrm{QBER} = \frac{N_{\text{errori}}}{N_{\text{sifted}}}
        \lesssim 11\% \ \text{(BB84)}
      \]
    \end{column}
    \begin{column}{0.42\textwidth}
      \begin{block}{Ciò che interessa alla rete}
        Il \emph{secret key rate} (SKR), in bit/s, cala con la distanza:
        \[
          \eta = 10^{-\alpha L/10}
        \]
        con \(\alpha \approx 0{,}2\)~dB/km.
      \end{block}
      \begin{alertblock}{Due osservazioni}
        \begin{itemize}
          \item La QKD \emph{distribuisce} chiavi, non cifra né autentica.
          \item Serve un segreto iniziale: è \emph{key growing}.
        \end{itemize}
      \end{alertblock}
    \end{column}
  \end{columns}
\end{frame}

\begin{frame}{Dal link alla rete: il trusted node}
  \begin{columns}[c]
    \begin{column}{0.52\textwidth}
      \begin{itemize}
        \item Un link copre al più qualche centinaio di km; non esistono
              ancora ripetitori quantistici pratici.
        \item Le reti reali collegano i link in serie tramite
              \alert{nodi fidati}.
        \item Ogni hop è un one-time pad, ma $K$ è in chiaro dentro ogni
              nodo intermedio.
        \item Il relay \alert{consuma chiave}: ogni link $\ell\in p$ spende
              $|K|$ bit. La chiave diventa una risorsa da pianificare.
      \end{itemize}
    \end{column}
    \begin{column}{0.46\textwidth}
      \begin{center}\begin{adjustbox}{max width=\linewidth}
      \begin{tikzpicture}[scale=0.85,
          nd/.style={draw,circle,fill=green!12,minimum size=10mm,font=\small}]
        \node[nd] (A) at (0,0) {$A$};
        \node[nd,fill=orange!20] (B) at (2.6,0) {$B$};
        \node[nd] (C) at (5.2,0) {$C$};
        \draw[thick] (A) -- node[above,font=\scriptsize]{$K_{AB}$} (B);
        \draw[thick] (B) -- node[above,font=\scriptsize]{$K_{BC}$} (C);
        \node[font=\scriptsize,align=center] at (1.3,-1.0) {$c = K\oplus K_{AB}$};
        \node[font=\scriptsize,align=center] at (3.9,-1.0) {$c' = K\oplus K_{BC}$};
        \node[font=\scriptsize,align=center,red!70!black] at (2.6,1.1)
             {$K = c\oplus K_{AB}$};
      \end{tikzpicture}
      \end{adjustbox}\end{center}
    \end{column}
  \end{columns}
\end{frame}

\subsection{Architettura e SDN}

\begin{frame}{L'architettura a strati ITU-T Y.3800}
  \slidefig[0.55\textheight]{itu_strati}
  \small
  \begin{columns}[T]
    \begin{column}{0.5\textwidth}
      \textbf{QKD link}: due moduli QKD, dove la chiave è \emph{generata}.\\
      \textbf{KM link}: due key manager, dove la chiave è \emph{inoltrata}.
    \end{column}
    \begin{column}{0.5\textwidth}
      \textbf{QKDN controller}: decide in tempo reale (cammini, link).\\
      \textbf{QKDN manager}: supervisiona (FCAPS, politiche).
    \end{column}
  \end{columns}
  \source{Adattata da ITU-T Y.3800~\cite{ituy3800}}
\end{frame}

\begin{frame}{Software-Defined Networking}
  \begin{columns}[c]
    \begin{column}{0.5\textwidth}
      \begin{itemize}
        \item SDN separa \emph{piano dati}, \emph{piano di controllo}
              (controller con vista globale) e \emph{piano applicativo}~\cite{kreutz2015}.
        \item Il controller parla ai dispositivi con la \alert{SBI} e offre
              servizi verso l'alto con la \alert{NBI}.
        \item Con più domini si aggiunge un \alert{orchestrator}.
      \end{itemize}
      \vspace{1ex}
      \small Una rete QKD ha tutto ciò che rende utile SDN: risorse scarse
      e variabili, vincoli fisici, apparati multi-vendor, domanda che cambia.
    \end{column}
    \begin{column}{0.46\textwidth}
      \begin{center}
      \begin{tikzpicture}[node distance=8mm]
        \node[box,fill=blue!8,minimum width=3.6cm] (o) {Orchestrator};
        \node[box,fill=green!8,minimum width=3.6cm,below=of o] (c) {Controller};
        \node[box,fill=gray!8,minimum width=3.6cm,below=of c] (d) {Dispositivi};
        \draw[<->,thick] (o) -- node[right,font=\scriptsize]{NBI\quad(QKD 018)} (c);
        \draw[<->,thick] (c) -- node[right,font=\scriptsize]{SBI\quad(QKD 015)} (d);
      \end{tikzpicture}
      \end{center}
    \end{column}
  \end{columns}
\end{frame}

% =====================================================================
\subsection{GS QKD 002}
% =====================================================================

\begin{frame}{GS QKD 002 (2010): dove serve la QKD}
  \begin{columns}[T]
    \begin{column}{0.48\textwidth}
      \begin{block}{Una primitiva tra tre}
        Un sistema sicuro richiede \emph{cifratura}, \emph{distribuzione
        delle chiavi}, \emph{autenticazione}. La QKD copre \alert{solo la
        seconda}: è un complemento, non un sostituto, della crittografia
        classica (OTP/AES, MAC Wegman–Carter).
      \end{block}
    \end{column}
    \begin{column}{0.5\textwidth}
      \textbf{Dove si consuma la chiave}
      \begin{description}\small
        \item[L2] PPP, MACsec, QKD link encryptor
        \item[L3] IKE modificato per IPsec
        \item[L4] chiavi di sessione TLS, HMAC
        \item[L7] autenticazione e cifratura applicativa
      \end{description}
      \small In ogni caso c'è chi \emph{consuma} e chi \emph{fornisce}
      chiavi: la futura separazione \alert{SAE / KME}.
    \end{column}
  \end{columns}
  \source{\cite{qkd002}}
\end{frame}

\begin{frame}{I sei casi d'uso}
  \scriptsize
  \begin{tabularx}{\textwidth}{@{}l L L@{}}
    \toprule
    \textbf{Caso} & \textbf{Configurazione} & \textbf{Esigenza che anticipa} \\
    \midrule
    UC1 Backup & link punto-punto + cifratore di linea & interfaccia chiave–cifratore \\
    UC2 Metro aziendale & più link, server di chiavi & QoS per applicazione \\
    UC3 Infrastrutture critiche & rete con key relay e gestione integrata & KMS, ridondanza, controllo centralizzato \\
    UC4 Backbone & canale quantistico in WDM & coesistenza con la rete ottica, gestione integrata \\
    UC5 Accesso PON & punto-multipunto & condivisione della capacità di chiave \\
    UC6 Long-haul & relay tramite nodo mobile (satellite) & trusted node, ridondanza \\
    \bottomrule
  \end{tabularx}
  \vspace{1ex}
  \footnotesize
  \begin{exampleblock}{Che cosa lascia in eredità}
    La QKD utile è una \emph{rete}: gestione delle chiavi (UC3, UC6) e QoS
    differenziata (UC2) $\Rightarrow$ \textbf{GS QKD 015}; integrazione con
    la rete ottica dell'operatore (UC4) $\Rightarrow$ \textbf{GS QKD 018}.
  \end{exampleblock}
\end{frame}

% =====================================================================
\section[GS QKD 015]{ETSI GS QKD 015: il controller SDN}
% =====================================================================

\subsection{Componenti}

\begin{frame}{Perché un controllo software-defined}
  \begin{columns}[T]
    \begin{column}{0.46\textwidth}
      \begin{alertblock}{Oggi: integrazione manuale}
        \begin{itemize}
          \item ogni sistema installato a mano su fibra scura dedicata;
          \item non scala;
          \item la capacità di chiave non si sposta dove serve;
          \item non dialoga con la rete ottica.
        \end{itemize}
      \end{alertblock}
    \end{column}
    \begin{column}{0.52\textwidth}
      \begin{block}{Obiettivi dell'astrazione}
        \begin{itemize}
          \item gestione centrale delle risorse QKD;
          \item un'unica gestione delle chiavi per perimetro;
          \item parametri di basso livello delegati al nodo;
          \item link distribuiti secondo la QoS;
          \item coordinamento di canali quantistici e classici.
        \end{itemize}
      \end{block}
    \end{column}
  \end{columns}
  \vspace{1ex}
  \centering\small GS QKD 015 V2.1.1 (2022)~\cite{qkd015}: la \alert{southbound
  interface} del piano di controllo QKD, allineata alla 018.
\end{frame}

\begin{frame}{Le componenti di una rete SD-QKD}
  \small
  \begin{columns}[T]
    \begin{column}{0.5\textwidth}
      \begin{description}[Interfaccia QKD]
        \item[Modulo QKD] hardware + software, specifico del vendor
        \item[Interfaccia QKD] \texttt{qkdi}: vista astratta di un modulo
              (tecnologia, modello, $\lambda$, attacco ottico)
        \item[Nodo SD-QKD] moduli nello stesso luogo sicuro, un solo
              \emph{SDN agent}, un key manager, un punto di accesso alle chiavi
      \end{description}
    \end{column}
    \begin{column}{0.5\textwidth}
      \begin{description}[Applicazione]
        \item[Link] \texttt{qkd\_link}: \texttt{PHYS} (canale quantistico
              diretto) o \texttt{VIRT} (key relay)
        \item[Applicazione] \texttt{qkd\_app}: \texttt{CLIENT} o
              \texttt{INTERNAL}, con QoS
        \item[Controller] SD-QNC: registro delle app, pianifica i link,
              dialoga con la rete ottica (SD-ONC)
      \end{description}
    \end{column}
  \end{columns}
  \vspace{1ex}
  \begin{alertblock}{Principio di sicurezza}
    Il materiale di chiave \emph{non lascia mai il perimetro del nodo}: al
    controller salgono solo identificativi, stati, tassi e parametri fisici.
  \end{alertblock}
\end{frame}

\begin{frame}{Il nodo SD-QKD}
  \slidefig[0.74\textheight]{nodo_sdqkd}
  \source{Adattata da GS QKD 015, Fig.~1~\cite{qkd015}}
\end{frame}

\begin{frame}{Link virtuali e applicazioni}
  \begin{columns}[T]
    \begin{column}{0.5\textwidth}
      \begin{block}{Key association link \texttt{VIRT}}
        \begin{itemize}
          \item consuma chiave dai link fisici sottostanti, quindi è
                un'\emph{applicazione interna} con
                \texttt{virt\_bandwidth};
          \item ogni nodo intermedio conosce
                \texttt{virt\_prev\_hop} e \texttt{virt\_next\_hop};
          \item richiede \texttt{key\_relay\_mode\_enable}.
        \end{itemize}
      \end{block}
    \end{column}
    \begin{column}{0.48\textwidth}
      \begin{block}{QoS di un'applicazione}
        \begin{itemize}
          \item \texttt{max\_bandwidth}, \texttt{min\_bandwidth} (bit/s)
          \item \texttt{jitter} massimo (ms)
          \item \texttt{ttl} della chiave nel key store
          \item condivisione di cammini o chiavi
        \end{itemize}
      \end{block}
      \small È la risposta all'esigenza di QoS differenziata di UC2.
    \end{column}
  \end{columns}
\end{frame}

\subsection{Modello YANG}

\begin{frame}{Anatomia di un modello YANG}
  YANG~\cite{rfc7950} è uno \alert{schema} (come XSD o JSON
      Schema) dei dati di configurazione e stato e delle operazioni.
  \begin{itemize}
        \item \texttt{leaf}, \texttt{leaf-list}, \texttt{container},
              \texttt{list} con \texttt{key}
        \item \texttt{config true} (decide il controller) /
              \texttt{config false} (misura il nodo)
        \item \texttt{grouping}/\texttt{uses} per riusare,
              \texttt{augment} per estendere un altro modulo
        \item \texttt{leafref}, \texttt{when}, \texttt{must} per vincoli
        \item \texttt{rpc} e \texttt{notification}
      \end{itemize}
\end{frame}

\begin{frame}[fragile]{YANG: intestazione del modulo e identity}
\begin{lstlisting}[language=yang]
module etsi-qkd-sdn-node {
  namespace "urn:etsi:qkd:yang:etsi-qkd-node";
  prefix "etsi-qkdn";
  import etsi-qkd-node-types { prefix "etsi-qkdn-types"; }
  revision "2022-01-30" { reference "ETSI GS QKD 015 V2.1.1"; }
  ...
}
identity QKD-LINK-TYPES { description "QKD key association link types."; }
identity VIRT { base QKD-LINK-TYPES; }
identity PHYS { base QKD-LINK-TYPES; }
typedef qkd-link-types {
  type identityref { base QKD-LINK-TYPES; }
}
\end{lstlisting}
  \footnotesize Le \texttt{identity} sono estendibili da altri moduli,
      a differenza di un'\texttt{enumeration}.
\end{frame}

\begin{frame}[fragile]{Il modello del nodo: \texttt{etsi-qkd-sdn-node}}
\begin{lstlisting}[basicstyle=\ttfamily\fontsize{6.8}{7.4}\selectfont]
module: etsi-qkd-sdn-node
  +--rw qkd_node
     +--rw qkdn_id              yang:uuid
     +--ro qkdn_status?         etsi-qkdn-types:qkdn-status-types
     +--rw qkdn_capabilities
     |  +--rw link_stats_support?  application_stats_support?
     |  +--rw key_relay_mode_enable?
     +--rw qkd_applications
     |  +--rw qkd_app* [app_id]
     |     +--rw app_id, app_type, server_app_id, client_app_id*
     |     +--rw app_qos            (max/min_bandwidth, jitter, ttl)
     |     +--rw backing_qkdl_id*, local_qkdn_id, remote_qkdn_id
     +--rw qkd_interfaces
     |  +--rw qkd_interface* [qkdi_id]
     |     +--rw qkdi_model, qkdi_type, qkdi_att_point
     |     +--rw qkdi_capabilities  (role, wavelength_range, ...)
     +--rw qkd_links
        +--rw qkd_link* [qkdl_id]
           +--rw qkdl_type, qkdl_enable, qkdl_local / qkdl_remote
           +--rw virt_prev_hop, virt_next_hop*, virt_bandwidth (VIRT)
           +--rw phys_channel_att, phys_wavelength          (PHYS)
           +--rw qkdl_performance (skr, eskr, expected_consumption)
\end{lstlisting}
\end{frame}

\begin{frame}{Il modello del nodo: sezioni e prestazioni}
      \small
      Un albero per nodo, quattro sezioni:
      \begin{itemize}
        \item identità e \textbf{capability};
        \item \textbf{applicazioni} con QoS;
        \item \textbf{interfacce} (moduli);
        \item \textbf{link} con prestazioni.
      \end{itemize}
      \begin{block}{Prestazioni del link $\ell$}
        \[
          \mathrm{EC}_\ell = \sum_{a\in\mathcal{A}_\ell} b_a
        \]
        \[
          \mathrm{ESKR}_\ell \approx \mathrm{SKR}_\ell - \mathrm{EC}_\ell
        \]
        {\footnotesize EC = \texttt{expected\_consumption}; lettura intuitiva,
        non formula dello standard.}
      \end{block}
\end{frame}

\begin{frame}{Notifiche dal nodo al controller}
  Interrogare continuamente i nodi è inefficiente: la 015 definisce
  notifiche YANG per oggetto.
  \vspace{1ex}
  \small
  \begin{tabularx}{\textwidth}{@{}l L@{}}
    \toprule
    \textbf{Oggetto} & \textbf{Notifiche} \\
    \midrule
    Applicazioni & \texttt{sdqkdn\_application\_new}, \texttt{\_qos\_update}, \texttt{\_disconnected} \\
    Interfacce & \texttt{sdqkdn\_interface\_new}, \texttt{\_down} (con \emph{reason}), \texttt{\_out} \\
    Link & \texttt{sdqkdn\_link\_down}, \texttt{\_perf\_update}, \texttt{\_overloaded} \\
    Generiche & eventi e allarmi estendibili \\
    \bottomrule
  \end{tabularx}
  \vspace{2ex}

  Inoltre NETCONF/RESTCONF permettono di sottoscrivere qualunque dato del
  datastore e ricevere ogni variazione in tempo reale.
\end{frame}

\subsection{Workflow e sicurezza}

\begin{frame}{Workflow: registrazione di un'applicazione}
  \slidefig[0.6\textheight]{seq_registrazione}
  \small Le applicazioni conoscono solo il proprio key manager: è il
  controller, \alert{registro centrale}, ad abbinare le richieste e creare
  l'\texttt{app\_id}.
  \source{Adattata da GS QKD 015, Fig.~2~\cite{qkd015}}
\end{frame}

\begin{frame}{Workflow: creazione di link fisici e virtuali}
  \begin{columns}[T]
    \begin{column}{0.48\textwidth}
      \begin{block}{Link fisico}
        \begin{enumerate}\small
          \item innesco (domanda sui nodi o NBI);
          \item due interfacce libere e compatibili;
          \item richiesta di un \emph{lightpath} all'SD-ONC;
          \item configurazione delle interfacce.
        \end{enumerate}
        \footnotesize Cuore della coesistenza con la rete ottica (UC4).
      \end{block}
    \end{column}
    \begin{column}{0.5\textwidth}
      \begin{block}{Link virtuale}
        \small Se manca un lightpath: cammino sulla rete QKD, nodi
        intermedi in modalità relay, estremi con la banda richiesta.
        Esempio (nostro, non dello standard):
        \[
          w_\ell = \frac{1}{\mathrm{ESKR}_\ell}
        \]
        \[
          p^\star = \arg\min_{p\in\mathcal{P}_r} \sum_{\ell\in p} w_\ell
        \]
        \[
          \mathcal{P}_r = \Big\{\, p : \min_{\ell\in p}\mathrm{ESKR}_\ell \ge r \,\Big\}
        \]
      \end{block}
    \end{column}
  \end{columns}
  \vspace{1ex}
  \small Link virtuali mal gestiti congestionano i link fisici: serve la
  vista globale del controller.
\end{frame}

\begin{frame}{Protocolli, sicurezza e limiti della 015}
  \begin{columns}[T]
    \begin{column}{0.48\textwidth}
      \textbf{Protocolli}
      \begin{itemize}\small
        \item NETCONF~\cite{rfc6241}: apparato–controller, transazionale,
              SSH/TLS, XML
        \item RESTCONF~\cite{rfc8040}: NBI, HTTP, CRUD, JSON/XML
      \end{itemize}
      \textbf{Installazione sicura}
      \begin{itemize}\small
        \item area sicura esistente: canale protetto con chiavi QKD;
        \item nuovo sito: segreto iniziale, poi chiavi QKD
              (\emph{key growing}).
      \end{itemize}
    \end{column}
    \begin{column}{0.5\textwidth}
      \begin{alertblock}{Che cosa resta fuori}
        La 015 controlla \emph{un} dominio con \emph{un} controller.
        \begin{itemize}\small
          \item più domini QKD, anche di vendor diversi;
          \item SAE nei nodi di un'altra rete (OTN);
          \item gestione con gli stessi strumenti di orchestrazione
                del resto della rete.
        \end{itemize}
      \end{alertblock}
      \centering$\Rightarrow$ \textbf{GS QKD 018}
    \end{column}
  \end{columns}
\end{frame}

% =====================================================================
\section[GS QKD 018]{ETSI GS QKD 018: l'interfaccia di orchestrazione}
% =====================================================================

\subsection{Scenario e modello}

\begin{frame}{GS QKD 018 V1.1.1 (2022)}
  \begin{itemize}
    \item Interfaccia tra \alert{SDN orchestrator} (client) e
          \alert{SDN controller} del dominio QKD (server)~\cite{qkd018}.
    \item È la \emph{northbound interface} del controller: il complemento
          della 015.
  \end{itemize}
  \[
    \text{orchestrator} \xrightarrow{\text{QKD 018}} \text{controller}
    \xrightarrow{\text{QKD 015}} \text{nodo SD-QKD}
    \xrightarrow{\text{QKD 004/014}} \text{SAE}
  \]
  \begin{columns}[T]
    \begin{column}{0.5\textwidth}
      \begin{block}{Orchestrazione SDN}
        \small Coordinamento automatico e continuo delle risorse per
        stabilire e rilasciare servizi end-to-end che attraversano più
        domini, ciascuno con il proprio controller.
      \end{block}
    \end{column}
    \begin{column}{0.5\textwidth}
      \begin{block}{Service link}
        \small Associazione logica di chiave tra due nodi QKD, single- o
        multi-hop: a livello di orchestrazione è il link virtuale della 015.
      \end{block}
    \end{column}
  \end{columns}
\end{frame}

\begin{frame}{Lo scenario di riferimento}
  \vspace{-1ex}
  \slidefig[0.46\textheight]{scenario_018}
  \vspace{-1.5ex}
  \small
  \begin{columns}[T]
    \begin{column}{0.5\textwidth}
      \begin{itemize}
        \item Rete QKD separata dalla rete di trasporto ottico (OTN):
              prestazioni, isolamento, responsabilità distinte.
        \item I \alert{SAE} stanno nei nodi OTN, le \alert{chiavi} nei
              nodi QKD; ogni controller vede solo la propria rete.
      \end{itemize}
    \end{column}
    \begin{column}{0.48\textwidth}
      L'orchestrator risolve:
      \begin{enumerate}
        \item \emph{address matching} nodo OTN $\leftrightarrow$ nodo QKD;
        \item provisioning end-to-end;
        \item gestione: topologia, inventario, FCAPS.
      \end{enumerate}
    \end{column}
  \end{columns}
  \source{Adattata da GS QKD 018, Fig.~1~\cite{qkd018}}
\end{frame}

\begin{frame}{Funzioni dell'orchestrator e modello informativo}
  \begin{columns}[T]
    \begin{column}{0.46\textwidth}
      \textbf{Funzioni richieste (§4.2)}
      \begin{itemize}\footnotesize
        \item tradurre le richieste in configurazione dei domini;
        \item attivare e rilasciare servizi inter-dominio;
        \item calcolare e modificare cammini multi-dominio;
        \item scoprire nodi e link, vista astratta;
        \item monitorare le risorse di chiave e FCAPS.
      \end{itemize}
    \end{column}
    \begin{column}{0.52\textwidth}
      \textbf{Cinque blocchi informativi (§5)}
      \begin{description}[§5.5]\footnotesize
        \item[§5.1] discovery della topologia, proattiva o reattiva
        \item[§5.2] stato e inventario: SKR, ESKR, ripristino
        \item[§5.3] stato dei servizi end-to-end
        \item[§5.4] provisioning con calcolo dei cammini candidati
        \item[§5.5] notifiche e soglie configurabili (QBER, SKR, ESKR)
      \end{description}
    \end{column}
  \end{columns}
  \vspace{1ex}
  \begin{block}{Ammissibilità di un cammino (formalizzazione nostra)}
    Ogni link di un cammino multi-hop inoltra la stessa chiave, quindi conta
    il collo di bottiglia:
    \[
      p \text{ ammissibile per la richiesta } r \iff
      \min_{\ell\in p} \mathrm{ESKR}_\ell \ge r
    \]
  \end{block}
\end{frame}

\subsection{Moduli YANG}

\begin{frame}[fragile]{Tre moduli YANG: topology}
  \begin{columns}[T]
    \begin{column}{0.54\textwidth}
\begin{lstlisting}[basicstyle=\ttfamily\fontsize{5.8}{6.6}\selectfont]
qkd_network
 +-- sdqkd_nodes
 |    +-- qkdn* [qkdn_id]               (nodi SD-QKD)
 |         +-- qkd_interfaces
 |              +-- qkdi* [qkdi_id]     (moduli QKD)
 +-- qkd_links
      +-- qkd_phys_links
      |    +-- phys_link* [phys_link_id]
      |         +-- link_type
      |         +-- phys_link_local_qkdn
      |         |    (qkdn_id_ref, qkdi_id_ref)
      |         +-- phys_link_remote_qkdn
      +-- qkd_svc_links
           +-- svc_link* [svc_link_id]
                +-- ... estremi locale/remoto
\end{lstlisting}
      \footnotesize Catena: \emph{topology} $\to$ \emph{inventory} $\to$
      \emph{connectivity}. I moduli importano i tipi della 015.
    \end{column}
    \begin{column}{0.44\textwidth}
      \small \texttt{etsi-qkd-sdn-topology}: un solo albero per l'intero
      dominio, cioè un grafo
      \[
        G_i = (V_i, E_i),\quad E_i = E_{\text{phys}}\cup E_{\text{svc}}
      \]
      con estremi \texttt{leafref} validati. L'orchestrator compone la vista
      multi-dominio:
      \[
        G = \Big(\bigcup_i V_i,\ \bigcup_i E_i \cup E_{\text{inter}}\Big)
      \]
      \footnotesize $E_{\text{inter}}$: adiacenze tra domini, tipicamente
      nodi di bordo.
    \end{column}
  \end{columns}
  \source{\cite{qkd018yang}}
\end{frame}

\begin{frame}{Tre moduli YANG: inventory}
  \small \texttt{etsi-qkd-sdn-inventory} non crea un nuovo albero: fa
  \alert{\texttt{augment}} della topologia.
  \vspace{1ex}

  \footnotesize
  \begin{tabularx}{\textwidth}{@{}l L@{}}
    \toprule
    \textbf{Elemento} & \textbf{Foglie aggiunte} \\
    \midrule
    \texttt{qkdn} & versione, capability, location, IP e porta del KME
    (dove i SAE prelevano le chiavi via QKD~014), stato \\
    \texttt{qkdi} & ruolo, modello, tipo, punto di attacco, stato \\
    \texttt{phys\_link} & \texttt{enable}, \texttt{wavelength},
    \texttt{performance} (\texttt{skr}, \texttt{eskr},
    \texttt{expected\_consumption}, QBER, SNR), stato \\
    \texttt{svc\_link} & \texttt{bandwidth}, \texttt{performance},
    \texttt{path\_restoration} (\texttt{AUTOMATIC}/\texttt{MANUAL},
    \texttt{hold\_off\_time}), stato \\
    \bottomrule
  \end{tabularx}
  \vspace{1.5ex}

  \small Notifiche: \texttt{sdqkdn\_new}, \texttt{\_interface\_new},
  \texttt{\_phys\_link\_new}, \texttt{\_svc\_link\_new}, varianti
  \texttt{\_perf\_update} e \texttt{\_overloaded}, \texttt{alarm},
  \texttt{event}.
\end{frame}

\begin{frame}[fragile]{Tre moduli YANG: connectivity}
  \begin{columns}[T]
    \begin{column}{0.52\textwidth}
\begin{lstlisting}[basicstyle=\ttfamily\fontsize{6}{6.8}\selectfont]
connectivity_service/qkd_service
 +-- qkd_app* [app_id]
      +-- (dati app da QKD 015: stato, QoS)
      +-- app_mapped_svc_link
           +-- svc_link_and_qkdn_ids
           +-- path_constraints
           |    (include/exclude_nodes)
           +-- candidate_paths
           +-- deploy_path (path_id)
\end{lstlisting}
    \end{column}
    \begin{column}{0.46\textwidth}
      \begin{block}{Divisione dei ruoli}
        \small Il controller conosce il dominio e \alert{propone} i
        cammini; l'orchestrator ha la vista d'insieme e \alert{decide}.
      \end{block}
    \end{column}
  \end{columns}
  \vspace{1ex}
  \textbf{RPC: il workflow di provisioning passo per passo}\par
      \vspace{0.5ex}
      \scriptsize
      \begin{tabular}{@{}r@{\ \ }l@{}}
        1 & \texttt{app-registration}\\
        2 & \texttt{create-service-link-id-for-app}\\
        3 & \texttt{create-candidate-paths-per-service-link-for-app}\\
        3' & \texttt{update-candidate-paths-per-service-link-for-app}\\
        4 & \texttt{deploy-a-candidate-path-per-service-link-for-app}\\
        5 & \texttt{delete-deployed-path-per-service-link-for-app}
      \end{tabular}
  \vspace{1ex}

  \textbf{Notifiche}\par
      \vspace{0.5ex}
      \scriptsize\texttt{sdqkdn\_application\_new}, \texttt{\_deployed},
      \texttt{\_qos\_update}, \texttt{\_disconnected}
\end{frame}

\begin{frame}{I tre moduli in sintesi}
  \small
  \begin{tabularx}{\textwidth}{@{}l L l L@{}}
    \toprule
    \textbf{Modulo} & \textbf{Domanda} & \textbf{Clausola} & \textbf{Dati} \\
    \midrule
    topology & Com'è fatta la rete? & §5.1 & struttura (grafo) \\
    inventory & Quante risorse ci sono e in che stato? & §5.2, §5.5 &
    stato, prestazioni, configurazione, notifiche \\
    connectivity & Come servo un'applicazione end-to-end? & §5.3, §5.4, §6.2 &
    applicazioni, RPC di calcolo e deploy dei cammini \\
    \bottomrule
  \end{tabularx}
  \vspace{3ex}
  \begin{exampleblock}{Perché questa struttura}
    Chi vuole solo la struttura implementa \emph{topology}; chi vuole anche
    lo stato aggiunge \emph{inventory} via \texttt{augment}, senza duplicare
    nulla; \emph{connectivity} aggiunge le operazioni.
  \end{exampleblock}
\end{frame}

\subsection{Workflow e sicurezza}

\begin{frame}{Workflow di provisioning QKD + OTN (§6.2)}
  \slidefig[0.76\textheight]{seq_018}
  \source{Rielaborazione semplificata di GS QKD 018, Fig.~2~\cite{qkd018}}
\end{frame}

\begin{frame}{Il workflow passo per passo}
  \begin{columns}[T]
    \begin{column}{0.56\textwidth}
      \begin{enumerate}\small
        \item L'amministratore chiede un servizio (SAE, QoS).
        \item L'orchestrator individua i controller ed effettua
              l'address matching.
        \item Chiede al controller QKD un service link e i cammini, con
              vincoli.
        \item Il controller assegna \texttt{svc\_link\_id} e restituisce
              candidati con priorità (passo ripetibile).
        \item L'orchestrator sceglie e chiede il deploy; i nodi fanno key
              relay lungo il cammino.
        \item Provisioning dell'app e avviso al controller OTN.
      \end{enumerate}
    \end{column}
    \begin{column}{0.42\textwidth}
      \begin{block}{015 vs 018}
        \small Nella 015 il controller calcola e configura da solo il link
        virtuale; nella 018 lo \alert{propone} e attende la scelta.
      \end{block}
      \begin{block}{Più domini QKD}
        \small L'orchestrator ripete i passi 3–5 per ogni dominio lungo il
        cammino inter-dominio, collegando i segmenti nei nodi di bordo.
      \end{block}
    \end{column}
  \end{columns}
\end{frame}

\begin{frame}[fragile]{Protocolli e sicurezza}
  \begin{columns}[T]
    \begin{column}{0.5\textwidth}
      \textbf{Trasporto (§8)}
      \begin{itemize}\small
        \item Modello (YANG 1.1) separato dal protocollo: NETCONF o
              RESTCONF.
        \item RESTCONF: ogni risorsa è un URI; \texttt{GET} per leggere,
              \texttt{POST}/\texttt{PUT}/\texttt{PATCH}/\texttt{DELETE}
              per configurare; notifiche via stream.
      \end{itemize}
    \end{column}
    \begin{column}{0.48\textwidth}
      \begin{alertblock}{Sicurezza (§7)}
        \small Le chiavi QKD \textbf{non transitano mai} per controller o
        orchestrator: nei modelli compaiono solo tassi e l'indirizzo del KME.
        Un orchestrator compromesso causa disservizi, non lettura delle chiavi.
        La protezione del canale (es.\ TLS) è lasciata all'implementazione.
      \end{alertblock}
    \end{column}
  \end{columns}
  \vspace{1ex}
\begin{lstlisting}
GET  /restconf/data/etsi-qkd-sdn-topology:qkd_network
POST /restconf/operations/etsi-qkd-sdn-connectivity:create-service-link-id-for-app
\end{lstlisting}
\end{frame}

% =====================================================================
\subsection{Analogia ACTN}
% =====================================================================

\begin{frame}{La 018 applica alla QKD il pattern ACTN}
  \footnotesize
  \begin{tabularx}{\textwidth}{@{}l l L@{}}
    \toprule
    \textbf{ACTN} & \textbf{ETSI QKD} & \textbf{Osservazione} \\
    \midrule
    CNC & amministratore / portale cliente & origina la richiesta \\
    MDSC & SDN orchestrator & vista multi-dominio, decisione del cammino \\
    PNC & SDN controller del dominio QKD & calcola i candidati, attiva i cammini \\
    MPI & interfaccia GS QKD 018 & topologia astratta, provisioning, notifiche \\
    SBI & interfaccia GS QKD 015 & controllo dei singoli nodi \\
    TE topology astratta & \texttt{etsi-qkd-sdn-topology} & vista vendor-neutral \\
    banda TE & SKR / ESKR & la risorsa da allocare è la chiave \\
    \bottomrule
  \end{tabularx}
  \vspace{1ex}
  \begin{alertblock}{La differenza: la risorsa}
    In ACTN la banda è fissata dalla capacità del canale; in QKD l'SKR
    \emph{fluttua} e la chiave si accumula in buffer finiti. Da qui il peso
    di \texttt{\_perf\_update}, \texttt{\_overloaded} e
    \texttt{path\_restoration}.
  \end{alertblock}
  \source{\cite{rfc8453}}
\end{frame}

% =====================================================================
\section[Casi reali e conclusioni]{Casi reali, limiti e conclusioni}
% =====================================================================

\subsection{SK Telecom}

\begin{frame}{SK Telecom: orchestrazione multi-vendor su KOREN}
  \begin{columns}[T]
    \begin{column}{0.52\textwidth}
      \begin{itemize}\small
        \item Primo dispiegamento documentato di un orchestrator
              conforme alla 018 (2020–2022)~\cite{sim2023}.
        \item Due operatori, KMS e dispositivi di vendor diversi, ognuno
              col proprio controller (QKDNC); un orchestrator (QKDNO).
        \item BB84 decoy state, tratte da 39 a 100~km.
        \item \alert{Nodi di bordo} (Pangyo, Daejeon): nessun modulo QKD,
              due KMS interconnessi.
      \end{itemize}
    \end{column}
    \begin{column}{0.46\textwidth}
      \begin{block}{Estensioni al modello}
        \small
        \begin{itemize}
          \item tipo di applicazione \emph{multi-segment} (\texttt{MS}),
                con \texttt{ms\_app\_id} comune ai domini;
          \item \texttt{server\_ms\_info}, \texttt{client\_ms\_info};
          \item RPC \texttt{app-registration-ms};
          \item chiavi al bordo via GS QKD 020~\cite{qkd020}.
        \end{itemize}
      \end{block}
      \footnotesize \texttt{CLIENT}/\texttt{INTERNAL} non bastano per il
      multi-dominio; le \texttt{identity} YANG permettono di estendere.
    \end{column}
  \end{columns}
\end{frame}

\subsection{PoliMi e MadQCI}

\begin{frame}{PoliMi e MadQCI}
  \begin{columns}[T]
    \begin{column}{0.49\textwidth}
      \begin{block}{Politecnico di Milano~\cite{vicario2025}}
        \small
        \begin{itemize}
          \item Controllo gerarchico per reti metropolitane.
          \item Controller di dominio su ONOS con il modello 015.
          \item Interfaccia 018 via RESTCONF verso un orchestrator Python
                che aggrega più domini.
          \item Tutta la pila 015 $\to$ 018 in software aperto.
        \end{itemize}
      \end{block}
    \end{column}
    \begin{column}{0.49\textwidth}
      \begin{block}{MadQCI, Madrid~\cite{madqci2024}}
        \small
        \begin{itemize}
          \item Rete reale: decine di moduli multi-vendor, oltre 130~km,
                integrata con la rete ottica di produzione.
          \item Controllo SD-QKD basato sulla 015.
          \item La 018 è indicata come attività \emph{ancora in corso}.
        \end{itemize}
      \end{block}
    \end{column}
  \end{columns}
  \vspace{2ex}
  \centering Il controllo di dominio è maturo, l'\alert{orchestrazione}
  molto meno.
\end{frame}

% =====================================================================
% =====================================================================

\subsection{Limiti}

\begin{frame}{Limiti e prospettive}
  \begin{columns}[T]
    \begin{column}{0.5\textwidth}
      \begin{description}[Multi-dominio]\small
        \item[Adozione] la 015 è in reti reali, la 018 quasi solo in
              testbed.
        \item[Multi-dominio] lo scenario normativo è QKD + OTN; QKD + QKD
              richiede estensioni e la GS QKD 020.
        \item[Metriche] l'ESKR è istantaneo e ignora i buffer~\cite{laschet2026}.
      \end{description}
    \end{column}
    \begin{column}{0.5\textwidth}
      \begin{description}[Riservatezza]\small
        \item[Riservatezza] esporre la topologia è accettabile dentro un
              operatore, meno tra operatori~\cite{laschet2026,interconnect2026}.
        \item[Sicurezza] protezione del canale orchestrator–controller non
              specificata.
        \item[Algoritmi] calcolo dei cammini non standardizzato~\cite{routing2024,mendez2026}.
      \end{description}
    \end{column}
  \end{columns}
\end{frame}

\subsection{Conclusioni}

\begin{frame}{Conclusioni}
  \begin{center}
  \begin{tikzpicture}[node distance=3mm,
      st/.style={draw,rounded corners,fill=blue!5,align=center,
                 font=\scriptsize,text width=2.9cm,minimum height=22mm}]
    \node[st] (a) {\textbf{GS QKD 002}\\[2pt] la QKD utile è una rete:
      gestione delle chiavi, QoS, integrazione ottica};
    \node[st,right=of a] (b) {\textbf{GS QKD 015}\\[2pt] il nodo SD-QKD:
      interfacce, key manager, link PHYS/VIRT, app con QoS, YANG};
    \node[st,right=of b,fill=orange!15] (c) {\textbf{GS QKD 018}\\[2pt]
      ogni dominio è un grafo astratto: il controller propone,
      l'orchestrator decide};
    \draw[->,thick] (a) -- (b);
    \draw[->,thick] (b) -- (c);
  \end{tikzpicture}
  \end{center}
  \vspace{1ex}
  \begin{itemize}
    \item Il disegno ricalca ACTN e riusa YANG, NETCONF e RESTCONF.
    \item Principio essenziale: \alert{le chiavi non lasciano mai i nodi}.
    \item KOREN, PoliMi e MadQCI: lo schema funziona, ma l'orchestrazione è
          giovane. Prossimo passo: multi-operatore, metriche migliori,
          riservatezza tra domini.
  \end{itemize}
\end{frame}

\begin{frame}[allowframebreaks]{Riferimenti}
  \printbibliography[heading=none]
\end{frame}

\begin{frame}
  \begin{center}
    {\Large Grazie per l'attenzione}\\[2ex]
    Domande?
  \end{center}
\end{frame}

\end{document}
